% !TEX program = xelatex

\documentclass{resume}
\usepackage{xeCJK}
\setCJKmainfont[BoldFont=STSongti-SC-Black]{STSongti-SC-Regular}
\setCJKsansfont{STHeitiSC-Medium}
\setlist{itemsep=0.15em, topsep=0.3em, leftmargin=1.5pc}
\setlist[itemize,2]{label=\textbullet}

% 第二、三页的项目块：标题 + 问题 / 行动 / 结果 / 改进
\newcommand{\starblock}[5]{%
  \par\vspace{0.55em}%
  {\large\textbf{#1}}\par\vspace{0.12em}%
  \textbf{要解决的问题：}#2\par
  \textbf{做了什么：}#3\par
  \textbf{结果：}#4\par
  \textbf{还可以怎么改进：}#5\par
}

\begin{document}
\pagenumbering{gobble}

\name{童浩}

\basicInfo{
  (+86) 188-5098-6208 \textperiodcentered\
  \href{mailto:1310554512@qq.com}{1310554512@qq.com} \textperiodcentered\
  中共党员
}

\section{个人摘要}
厦门理工学院计算机技术硕士，福州大学计算机本科。商汤科技算法研究实习，负责表格、笔记和 PPT 的大模型落地。另独立交付 AI 客服，覆盖知识库、多轮对话、实时语音和人工接管。先对齐业务问题，再用自动评测和人工审核稳住结果；开发时使用模型生成、对拍校验和回归测试。

\section{我会什么}
\begin{itemize}
  \item \textbf{业务：}智能客服、知识库问答、办公文档（表格 / 笔记 / PPT）、多智能体。
  \item \textbf{基础：}Python、JavaScript、SQL、C、C++；数据生成与标注、向量与关键词检索、接口与数据库、评测与回归。
  \item \textbf{用 AI 开发：}用模型生成样本和改写草稿，多次生成对拍，自动评分加人工抽检，改配置先看差异再发布，评测和功能测试分开。
\end{itemize}

\section{项目一览}
\datedsubsection{\textbf{商汤科技}，算法研究实习生}{2025.06 -- 至今}
\begin{itemize}
  \item \textbf{表格数据与标注：}补齐种子数据集中全部表格类型、分布均衡的训练样本，并用标注网站做质量把关。
  \item \textbf{表格评测与结果校验：}自动评测与人工评分一致性 92.4\%，日均节省评测人力超过 4 小时；推理结果验证准确率由 85.7\% 提升至 97.9\%。
  \item \textbf{生成质量：}字段命中率由 70.3\% 提升至 89.6\%，整体生成准确率提升 8.4\%。
  \item \textbf{笔记：}搜索回答带编号引用，能指出来源；导入主流笔记时，内容和元信息保持完整。
  \item \textbf{PPT：}两分钟内生成大纲和幻灯片结构，并支持套用用户自己的模板。
\end{itemize}

\datedsubsection{\textbf{AI 客服中台}，独立开发}{2026.08 -- 至今}
\begin{itemize}
  \item 面向正式交付，覆盖知识库问答、多轮文本、实时语音、人工接管和网站嵌入。
\end{itemize}

\section{教育与奖项}
\datedsubsection{\textbf{厦门理工学院}，计算机技术，工学硕士}{2023.09 -- 2026.06}
\datedsubsection{\textbf{福州大学}，计算机科学与技术，工学学士}{2018.09 -- 2022.06}
排名 20/151（前 15\%）。福州大学数学与计算机科学学院学生会副主席、校学生会网络科技部部长；国家励志奖学金、一等奖学金、三好学生、优秀共青团员。

第五届全国大学生集成电路创新创业大赛决赛\textbf{全国二等奖}（2021）；第十一届中国大学生服务外包创新创业大赛\textbf{全国三等奖}（2020）。

\newpage

\section{项目经历}
\datedsubsection{\textbf{商汤科技（SenseTime）}}{2025.06 -- 至今}
\role{项目角色}{算法研究实习生}

\starblock{表格数据与标注}{%
  表格模型要有类型齐全、分布均衡的训练数据，标注结果还要经得起复查。%
}{%
  以种子数据集为基准，用大语言模型和规则模板脚本生成多样化表格，覆盖种子集全部表格类型并保持分布均衡。标注网站对同一内容采样 K 次做 AI 对拍，比对后筛出异常再复审，并支持多角色协作。%
}{%
  训练样本覆盖种子数据集的全部表格类型，分布保持均衡；标注异常可以自动筛出并进入复审。%
}{%
  下一步按难度给样本分层，并记下抽检比例，用来区分简单表和难例。%
}

\starblock{表格评测与推理结果校验}{%
  人工评模型耗时长、打分不一致；表格推理结果遇到格式扰动时也不够稳。%
}{%
  按结构化维度给模型输出打分，加入人工评分修正后形成加权结果，用于 A/B 验证和回归测试。推理结果按 JAS（JavaScript for WPS Application）执行输出做哈希比对和语义比对。%
}{%
  系统评测与人工评分一致性达到 92.4\%，日均节省评测人力超过 4 小时。验证准确率由 85.7\% 提升至 97.9\%，对格式扰动更稳。%
}{%
  下一步按业务坏例补评测维度，确认高分结果在真实表格任务里也能用。%
}

\starblock{生成质量（GRPO 奖励）}{%
  模型输出经常缺少业务要求的结构化字段。%
}{%
  在 GRPO 训练中加入基于正则的分层级关键词奖励，引导模型写出必要的结构化信息。%
}{%
  字段命中率从 70.3\% 提升至 89.6\%，整体生成准确率提升 8.4\%。结构缺失和格式错误下降，\mbox{bad case} 总体占比减少 10\%，收敛更快。%
}{%
  下一步在关键词命中之外加入语义等价判定，使意思正确但表述不同的输出不被误罚。%
}

\starblock{笔记的搜索与格式转换}{%
  搜索回答需要能看出来源；从其他笔记应用导入时，正文和元信息要一起留下来。%
}{%
  搜索按 Perplexity 架构重写查询，用关键词和语义重排后取 top K 片段，生成带编号引用的回答。格式转换把模型结果映射为前端 HTML；导入时解析多种笔记格式，统一成内部结构。%
}{%
  回答可以回溯到带编号的来源片段；笔记内容和元信息能够完整迁移。%
}{%
  下一步补充搜索采纳率和引用点击，用来确认可追溯回答是否被实际使用。%
}

\starblock{PPT 多智能体}{%
  需要自动生成 PPT 大纲和每页结构，并套用用户自己的模板。%
}{%
  用多智能体工作流生成大纲和各页结构化信息。用户上传模板后，识别其中的版式和嵌套逻辑，生成对应模板文件，规则覆盖目前已知的幻灯片类型。%
}{%
  整体响应控制在 2 分钟以内，模板规则覆盖目前已知的幻灯片类型。%
}{%
  下一步补充用户修改率和版式采纳情况，用来确认生成结果能否直接使用。%
}

\newpage

\datedsubsection{\textbf{AI 客服中台}}{2026.08 -- 至今}
\role{项目角色}{独立开发}
\href{https://github.com/EmiyaredA/AI-customer}{github.com/EmiyaredA/AI-customer}

\textbf{技术栈：}Python、FastAPI、PostgreSQL、Chroma、React、TypeScript、WebSocket。

面向正式交付的 AI 客服。Studio、嵌入页和浏览器扩展共用同一套运行约定；会话协调以 PostgreSQL 为准，Chroma 和上传文件可以按发布结果重建。

\starblock{知识库问答}{%
  客服回答要跟已发布的知识一致，不能脱离原文自由发挥，也不能每次把全文塞进上下文。%
}{%
  分块与向量化分开，支持结构切分、问答抽取和模型索引摘要。运行时不预取全文，由 Agent 用多条检索词查知识库；单库内融合向量检索和 BM25，再按发布快照去重截断。生成式检索以完整正文为事实源，模型只选择证据范围，浏览目录、搜索片段和读取原文由确定性工具完成。%
}{%
  检索评测与功能测试分开，生成式库和向量库使用相同检索词成对比较。%
}{%
  下一步补充线上解决率和答非所问比例，用来判断知识库回答是否解决了用户问题。%
}

\starblock{配置修改与发布}{%
  调整机器人配置和知识库时，未审核的改动不能直接覆盖正在服务的版本。%
}{%
  根据真实会话整理证据并生成改进草稿，可限定修改字段、引用重点会话，并排对比后由人审核。写入和发布以版本号一致为前提，冲突时返回当前快照且不自动覆盖。渠道发布在同一事务中提交版本、知识关联和运行数据，线上只读取不可变快照。配置以 BotBuildSpec 为唯一来源，由 OpenAPI 生成管理端类型，Studio、嵌入页和扩展共用同一运行时客户端。%
}{%
  版本冲突不会自动覆盖线上配置，运行时只读取已发布快照。%
}{%
  下一步用同一批会话对比发布前后的回答，确认配置改动是否提高解决率。%
}

\starblock{人工接管、实时语音与网站嵌入}{%
  忙时人工要能接手当前会话，语音不能误打断用户，网站嵌入不能把渠道密钥放在页面里。%
}{%
  REST 与 SSE 共用回合协调：先保存用户消息，模型和工具执行结束后再一次性写入助手回复、状态和日志；断线或取消只停止尚未开始的检查点。坐席接管后该会话停用 AI，邀约、心跳和恢复由持久任务推进，结案后按原版本恢复。试用页和嵌入页共用浏览器插话检测，DashScope 与 SenseAudio 的协议差异放在后端适配层；本地语音活动检测先做可恢复暂停，疑似回声从原处继续播放，只有暂停后仍判定为近端人声才通知供应商打断。嵌入由接入方后端保存渠道密钥，页面只换取短时 token。%
}{%
  语音测试集包含 80 个语义脚本、139 条干净音频和 200 个分层场景，用于检验插话和识别，不是线上业务效果。%
}{%
  下一步补充转人工率和误打断比例，用来确认接管和插话是否达到上线要求。%
}

\end{document}
